\begin{figure*}[t]
    \centering
    \subfigure[Transfer from PTB to Genia.]{\label{fig:genia.ptb}\includegraphics[width=42mm]{figures/genia_ptb_txt.png}} ~~~
    \subfigure[Transfer from CoNLL 2003 NER to Genia.]{\label{fig:genia.ner}\includegraphics[width=42mm]{figures/genia_ner_txt.png}} ~~~
    \subfigure[Transfer from Spanish NER to Genia.]{\label{fig:genia.span}\includegraphics[width=42mm]{figures/genia_span_txt.png}} \\
    \subfigure[Transfer from PTB to Twitter POS tagging.]{\label{fig:twitter.pos}\includegraphics[width=42mm]{figures/twitter_pos_txt.png}} ~~~
    \subfigure[Transfer from CoNLL 2003 to Twitter NER.]{\label{fig:twitter.ner}\includegraphics[width=42mm]{figures/twitter_ner_txt.png}} \\
    \subfigure[Transfer from CoNLL 2003 NER to PTB POS tagging.]{\label{fig:pos.ner}\includegraphics[width=42mm]{figures/pos_ner_txt.png}} ~~~
    \subfigure[Transfer from PTB POS tagging to CoNLL 2000 chunking.]{\label{fig:chunk.pos}\includegraphics[width=42mm]{figures/chunk_pos_txt.png}} ~~~
    \subfigure[Transfer from PTB POS tagging to CoNLL 2003 NER.]{\label{fig:ner.pos}\includegraphics[width=42mm]{figures/ner_pos_txt.png}} \\
    \subfigure[Transfer from CoNLL 2003 English NER to Spanish NER.]{\label{fig:span.eng}\includegraphics[width=42mm]{figures/span_eng_txt.png}} ~~~
    \subfigure[Transfer from Spanish NER to CoNLL 2003 English NER.]{\label{fig:eng.span}\includegraphics[width=42mm]{figures/eng_span_txt.png}}
    \caption{\small Results on transfer learning. Cross-domain transfer: Figures \ref{fig:genia.ptb}, \ref{fig:twitter.pos}, and \ref{fig:twitter.ner}. Cross-application transfer: Figures \ref{fig:pos.ner}, \ref{fig:chunk.pos}, and \ref{fig:ner.pos}. Cross-lingual transfer: Figures \ref{fig:span.eng} and \ref{fig:eng.span}. Transfer across domains and applications: Figure \ref{fig:genia.ner}. Transfer across domains, applications, and languages: Figure \ref{fig:genia.span}.}\label{fig:perf}
\end{figure*}

\begin{table*}[t]
\vspace{-0.1in}
\caption{\small Dataset statistics.}
\label{tab:stat}
\begin{center}
\begin{tabular}{lllrrr}
Benchmark & Task & Language & \# Training Tokens & \# Dev Tokens & \# Test Tokens
\\ \hline \\
PTB 2003 & POS Tagging & English &  912,344 & 131,768 & 129,654 \\
CoNLL 2000 & Chunking & English & 211,727 & - & 47,377 \\
CoNLL 2003 & NER & English & 204,567 & 51,578 & 46,666 \\
CoNLL 2002 & NER & Dutch & 202,931 & 37,761 & 68,994 \\
CoNLL 2002 & NER & Spanish & 207,484 & 51,645 & 52,098 \\
Genia & POS Tagging & English & 400,658 & 50,525 & 49,761 \\
Twitter & POS Tagging & English & 12,196 & 1,362 & 1,627 \\
Twitter & NER & English & 36,936 & 4,612 & 4,921 \\
\end{tabular}
\end{center}
\end{table*}

\begin{table*}[t]
\vspace{-0.1in}
\caption{\small Improvements with transfer learning under multiple low-resource settings (\%). ``Dom'', ``app'', and ``ling'' denote cross-domain, cross-application, and cross-lingual transfer settings respectively. The numbers following the slashes are labeling rates (chosen such that the number of labeled examples are of the same scale).}
\label{tab:imp}
\begin{center}
\begin{tabular}{lllllll}
Source & Target & Model & Setting & Transfer & No Transfer & Delta
\\ \hline \\
PTB & Twitter/0.1 & T-A & dom & 83.65 & 74.80 & 8.85 \\
CoNLL03 & Twitter/0.1 & T-A & dom & 43.24 & 34.65 & 8.59 \\
PTB & CoNLL03/0.01 & T-B & app & 74.92 & 68.64 & 6.28 \\
PTB & CoNLL00/0.01 & T-B & app & 86.73 & 83.49 & 3.24 \\
CoNLL03 & PTB/0.001 & T-B & app & 87.47 & 84.16 & 3.31 \\
Spanish & CoNLL03/0.01 & T-C & ling & 72.61 & 68.64 & 3.97 \\
CoNLL03 & Spanish/0.01 & T-C & ling & 60.43 & 59.84 & 0.59 \\
\hline \\
PTB & Genia/0.001 & T-A & dom & 92.62 & 83.26 & 9.36 \\
CoNLL03 & Genia/0.001 & T-B & dom\&app & 87.47 & 83.26 & 4.21 \\
Spanish & Genia/0.001 & T-C & dom\&app\&ling & 84.39 & 83.26 & 1.13 \\
PTB & Genia/0.001 & T-B & dom & 89.77 & 83.26 & 6.51 \\
PTB & Genia/0.001 & T-C & dom & 84.65 & 83.26 & 1.39 \\
\end{tabular}
\end{center}
\end{table*}


\begin{table*}[t]
\caption{\small Comparison with state-of-the-art results (\%).}
\label{tab:sota}
\begin{center}
\begin{tabular}{llllll}
Model & CoNLL 2000 & CoNLL 2003 & Spanish & Dutch & PTB 2003
\\ \hline \\
% Lin and Wu \shortcite{lin2009phrase} & -- & 90.90 & -- & -- & -- \\
\cite{collobert2011natural} & 94.32 & 89.59 & -- & -- & 97.29 \\
\cite{passos2014lexicon} & -- & 90.90 & -- & -- & -- \\
\cite{luo2015joint} & -- & 91.2 & -- & -- & -- \\
\cite{huang2015bidirectional} & 94.46 & 90.10 & -- & -- & 97.55 \\
\cite{gillick2015multilingual} & -- & 86.50 & 82.95 & 82.84 & -- \\
\cite{ling2015finding} & -- & -- & -- & -- & \textbf{97.78} \\
\cite{lample2016neural} & -- & 90.94 & 85.75 & 81.74 & -- \\
\cite{ma2016end} & -- & 91.21 & -- & -- & 97.55 \\
\hline
Ours w/o transfer & 94.66 & 91.20 & 84.69 & 85.00 & 97.55 \\
Ours w/ transfer & \textbf{95.41} & \textbf{91.26} & \textbf{85.77} & \textbf{85.19} & 97.55 \\
\end{tabular}
\end{center}
\vspace{-0.1in}
\end{table*}

\section{Experiments}

\subsection{Datasets}

We use the following benchmark datasets in our experiments: Penn Treebank (PTB) POS tagging, CoNLL 2000 chunking, CoNLL 2003 English NER, CoNLL 2002 Dutch NER, CoNLL 2002 Spanish NER, the Genia biomedical corpus \citep{kim2003genia}, and a Twitter corpus \citep{ritter2011named}. The statistics of the datasets are described in Table \ref{tab:stat}.
We construct the POS tagging dataset with the instructions described in \cite{toutanova2003feature}. Note that as a standard practice, the POS tags are extracted from the parsed trees. For the CoNLL 2003 English NER dataset, we follow previous works \citep{collobert2011natural} to append one-hot gazetteer features to the input of the CRF layer for fair comparison. 
% We truncate all words whose character sequence length is longer than a threshold (17 for English, 35 for Dutch, 20 for Spanish, and 35 for Genia and Twitter). We replace all numeric characters with ``0''. We use the BIOES (Begin, Inside, Outside, End, Single) tagging scheme \citep{ratinov2009design}. 
Since there is no standard training/dev/test data split for the Genia and Twitter corpora, we randomly sample 10\% for test, 10\% for development, and 80\% for training. We follow previous work \citep{barrett2014token} to map Genia POS tags to PTB POS tags.

\subsection{Transfer Learning Performance} \label{sec:transfer}

We evaluate our transfer learning approach on the above datasets. We fix the hyperparameters for all the results reported in this section: we set the character embedding dimension at $25$, the word embedding dimension at $50$ for English and $64$ for Spanish, the dimension of hidden states of the character-level GRUs at $80$, the dimension of hidden states of the word-level GRUs at $300$, and the initial learning rate at $0.01$. Except for the Twitter datasets, these datasets are fairly large.  To simulate a low-resource setting, we also use random subsets of the data. We vary the labeling rate of the target task at $0.001$, $0.01$, $0.1$ and $1.0$. Given a labeling rate $r$, we randomly sample a ratio $r$ of the sentences from the training set and discard the rest of the training data---e.g., a labeling rate of $0.001$ results in around 900 training tokens on PTB POS tagging (Cf. Table \ref{tab:stat}).

The results on transfer learning are plotted in Figure \ref{fig:perf}, where we compare the results with and without transfer learning under various labeling rates. The numbers in the y-axes are accuracies for POS tagging, and chunk-level F1 scores for chunking and NER. The numbers are shown in Table \ref{tab:imp}. We can see that our transfer learning approach consistently improved over the non-transfer results. We also observe that the improvement by transfer learning is more substantial when the labeling rate is lower. For cross-domain transfer, we obtained substantial improvement on the Genia and Twitter corpora by transferring the knowledge from PTB POS tagging and CoNLL 2003 NER. For example, as shown in Figure \ref{fig:genia.ptb}, we can obtain an tagging accuracy of $83\%+$ with zero labels and $92\%$ with only $0.001$ labels when transferring from PTB to Genia. As shown in Figures \ref{fig:twitter.pos} and \ref{fig:twitter.ner}, our transfer learning approach can improve the performance on Twitter POS tagging and NER for all labeling rates, and the improvements with $0.1$ labels are more than $8\%$ for both datasets. Cross-application transfer also leads to substantial improvement under low-resource conditions. For example, as shown in Figures \ref{fig:chunk.pos} and \ref{fig:ner.pos}, the improvements with $0.1$ labels are $6\%$ and $3\%$ on CoNLL 2000 chunking and CoNLL 2003 NER respectively when transferring from PTB POS tagging. Figures \ref{fig:eng.span} and \ref{fig:span.eng} show that cross-lingual transfer can improve the performance when few labels are available.

Figure \ref{fig:perf} further shows that the improvements by different architectures are in the following order: T-A $>$ T-B $>$ T-C. This phenomenon can be explained by the fact that T-A shares the most model parameters while T-C shares the least. 
% Therefore T-A can improve the most since the most knowledge (model parameters and feature representations) can be transferred.
Transfer settings like cross-lingual transfer can only use T-C because the underlying similarities between the source task and the target task are less prominent (i.e., less \textit{transferable}), and in those cases the improvement by transfer learning is less substantial.

Another interesting comparison is among Figures \ref{fig:genia.ptb}, \ref{fig:genia.ner}, and \ref{fig:genia.span}. Figure \ref{fig:genia.ptb} is cross-domain transfer, Figure \ref{fig:genia.ner} is transfer across domains and applications at the same time, and Figure \ref{fig:genia.span} combines all the three transfer settings (i.e., from Spanish NER in the general domain to English POS tagging in the biomedical domain). The results show that the improvement by transfer learning diminishes when the transfer becomes ``indirect'' (i.e., the source task and the target task are more loosely related).

We also study using different transfer learning models for the same task. We study the effects of using T-A, T-B, and T-C when transferring from PTB to Genia, and the results are included in the lower part of Table \ref{tab:imp}. We observe that the performance gain decreases when less parameters are shared (i.e., T-A $>$ T-B $>$ T-C).

\subsection{Comparison with State-of-the-Art Results}

In the above section, we examine the effects of different transfer learning architectures. Now we compare our approach with state-of-the-art systems on these datasets.

We use publicly available pretrained word embeddings as initialization. On the English datasets, following previous works that are based on neural networks \citep{collobert2011natural,huang2015bidirectional,chiu2015named,ma2016end}, we experiment with both the 50-dimensional SENNA embeddings \citep{collobert2011natural} and the 100-dimensional GloVe embeddings \citep{pennington2014glove} and use the development set to choose the embeddings for different tasks and settings. For Spanish and Dutch, we use the 64-dimensional Polyglot embeddings \citep{al2013polyglot}. We set the hidden state dimensions to be 300 for the word-level GRU. The initial learning rate for AdaGrad is fixed at 0.01. We use the development set to tune the other hyperparameters of our model. 

Our results are reported in Table \ref{tab:sota}. Since there are no standard data splits on the Genia and Twitter corpora, we do not include these datasets into our comparison. The results for CoNLL 2000 chunking, CoNLL 2003 NER, and PTB POS tagging are obtained by transfer learning between the three tasks, i.e., transferring from two tasks to the other. The results for Spanish and Dutch NER are obtained with transfer learning between the NER datasets in three languages (English, Spanish, and Dutch).
From Table \ref{tab:sota}, we can draw two conclusions. First, our transfer learning approach achieves new state-of-the-art results on all the considered benchmark datasets except PTB POS tagging, which indicates that transfer learning can still improve the performance even on datasets with relatively abundant labels. Second, our base model (w/o transfer) performs competitively compared to the state-of-the-art systems, which means that the improvements shown in Section \ref{sec:transfer} are obtained over a strong baseline.



